% =====================================================================================
% macros.tex -- notation and theorem environments of the article.
% Loaded by main.tex after the packages.  Each quantity listed here has one symbol,
% used throughout the article.
%
% NOTATION TABLE
% -------------------------------------------------------------------------------------
%  symbol (macro)                meaning
% -------------------------------------------------------------------------------------
%  \Om, \dOm                     spatial domain Omega and its boundary
%  \QT                           space-time cylinder (0,T] x Omega (evolution problems)
%  d                             spatial dimension (Section 6: d = 2,...,20; cost up to 100;
%                                constants and bias up to d = 10^6)
%  \x                            point of Omega (bold x); x_i its coordinates
%  \lap, \grad, \dn              Laplacian, gradient, outward normal derivative
%  \uex                          exact solution u^*
%  \unet                         network approximation u_theta (theta = all weights)
%  \ubeta                        minimiser of the penalised Ritz energy with weight beta
%  \Lpinn                        PINN loss (empirical, eq. s3_methods:eq:pinn)
%  \Lpop                         population (continuum) PINN loss
%  \Eritz                        penalised Ritz energy E_beta (eq. s3_methods:eq:ritz)
%  \lam                          weight of the boundary term in the PINN loss
%                                (lambda = 1 in Sections 4-5; tuned in Section 6)
%  \bet                          boundary-penalty weight beta of the Deep Ritz energy
%  \Nr, \Nb, \Nic                numbers of interior (residual), boundary and
%                                initial-condition collocation points
%  \errL                         relative L2 error  ||u_theta-u^*||_2 / ||u^*||_2
%                                on the held-out evaluation set
%  \errI                         L-infinity error  max |u_theta-u^*| on the evaluation set
%  \errC                         centred relative L2 error ||u_theta-u^*|| / ||u^*-mean(u^*)||
%  \errH                         relative H1-seminorm error ||grad(u_theta-u^*)|| / ||grad u^*||
%  \errT(t)                      relative L2 error on the time slice t (heat problem)
%  \norm{.}, \abs{.}, \ip{.}{.}  norm, absolute value, L2 inner product
%  \R, \N                        real numbers, natural numbers
%  \dd                           upright differential d
%  \ee                           Euler's number (upright e)
%  \sci{a}{b}                    a x 10^b, e.g. \sci{1.6}{-3}
%  \PH, \PWone, \PWtwo,          problem names: Heat (1-D heat), Wave1 (two-mode 1-D wave),
%  \PLtwo, \PLthree,             Wave2 (2-D wave), Lap2 (2-D Laplace, mixed BC), Lap3 (3-D
%  \PLthreeS, \Pone, \Ptwo       Laplace, discontinuous data), Lap3s (3-D Laplace, compatible
%                                data), LapD (d-dim Laplace), PoiD (d-dim Poisson).
%  \Pridge, \Pcospair,          further problems on the unit cube of Section 6.8 (RidgeD,
%  \Paltridge, \Psinlin,         CosPairD, AltRidgeD; SinLinD, ExpCosD of the re-check), held
%  \Pexpcos                      out from the design of the two remedies tested there
%                                The digit is the space dimension.  The names are words so
%                                that they cannot be confused with the norm L^2, the space
%                                H^1 or the polynomial spaces P_0, P_1.
%  \strat{name}                  collocation strategy name, one of
%                                grid-c (cell-centred grid), grid-n (node-centred, closed grid),
%                                random (fixed i.i.d. uniform), resample (fresh i.i.d. uniform
%                                every Adam step), Sobol (scrambled Sobol', fixed),
%                                RAD (residual-based adaptive distribution)
%  \Adam, \LBFGS, \AdamLBFGS     optimiser arms: Adam; L-BFGS; Adam followed by L-BFGS
%  \nbd{.}                       norm of L2(dOmega) (surface measure)
%  \kappa_d                      norm of the harmonic extension L2(dOmega) -> L2(Omega) on the unit
%                                cube (Section 6.7); not the constant C_d = sqrt(1+1/(d pi^2)) of
%                                Proposition s3_methods:prop:robin(c)
%  U_d, L_d                      upper and lower bounds of kappa_d^2 (Theorem s6_highdim:thm:kappa)
%  \delta_1                      first biharmonic Steklov eigenvalue of the unit cube (= 1/kappa_d^2)
%  B_int, B_bd, \eta             interior and boundary terms of the computable error bound and its
%                                efficiency (bound over error), Corollary s6_highdim:cor:bound(c);
%                                eta is used for nothing else
%  \Lpop_int, \Lpop_bd           the two terms of the population PINN loss \Lpop_1 (mean squared residual,
%                                mean squared boundary misfit), Corollary s6_highdim:cor:bound(c)
%  m_d                           bound of kappa_d^2 through the torsion function \vartheta of the cube
%                                (-Lap theta = 1, theta = 0 on the boundary), Remark sA_proofs:rem:steklov
%  q, N                          q = d_n u^* on the boundary and N = ||q||_partial (Section S7.5)
%  M_\partial                    boundary Gram matrix of the affine functions (Section S7.6)
%  \mathcal A, \mathcal F, \mathcal S   affine functions, span of the lift features, additively separable
%                                functions (Section S7.6); sinc a = sin(a)/a
%  \rho                          seed-paired ratio err(plain)/err(arm) of Section 6.8
%  "plateau"                     (Section 6.8) a test error no smaller than the affine floor
%  "the re-check"                (Sections 6.7, 6.8, S5, S7) the two rounds of re-checks of the analyses of
%                                Sections 6.7 and 6.8, made with separately written code (Section 3.9)
%  \code{...}                    file names, identifiers (typewriter)
% -------------------------------------------------------------------------------------
% Terminology fixed for the whole article:
%   "iteration"   = one optimiser step (never "epoch");
%   "network"     = one trained network up to the point where the optimiser arms branch,
%                   i.e. one (problem, strategy, seed); "run" = one network in one arm;
%                   "pair" = the two runs of one network;
%   "held-out"    = evaluation points never used for training;
%   "seed"        = fixes initial weights and all random point sets of a run;
%   "relative L2 error" always means \errL on the held-out set unless a slice is named.
% =====================================================================================

% ---- sets, operators -------------------------------------------------------------
\newcommand{\R}{\mathbb{R}}
\newcommand{\N}{\mathbb{N}}
\newcommand{\dd}{\mathrm{d}}
\newcommand{\ee}{\mathrm{e}}
\newcommand{\Om}{\Omega}
\newcommand{\dOm}{\partial\Omega}
\newcommand{\QT}{Q_T}
\newcommand{\x}{\boldsymbol{x}}
\newcommand{\lap}{\Delta}
\newcommand{\grad}{\nabla}
\newcommand{\dn}{\partial_n}
\newcommand{\norm}[1]{\lVert #1\rVert}
\newcommand{\abs}[1]{\lvert #1\rvert}
\newcommand{\ip}[2]{\langle #1,\,#2\rangle}
\newcommand{\nbd}[1]{\norm{#1}_{\partial}}

% ---- solutions and functionals ---------------------------------------------------
\newcommand{\uex}{u^{*}}
\newcommand{\unet}{u_{\theta}}
\newcommand{\ubeta}{u_{\beta}}
\newcommand{\Lpinn}{\mathcal{L}}
\newcommand{\Lpop}{\mathcal{L}^{\mathrm{pop}}}
\newcommand{\Eritz}{\mathcal{E}_{\beta}}
\newcommand{\lam}{\lambda}
\newcommand{\bet}{\beta}

% ---- point budgets ---------------------------------------------------------------
\newcommand{\Nr}{N_{r}}
\newcommand{\Nb}{N_{b}}
\newcommand{\Nic}{N_{0}}

% ---- error measures --------------------------------------------------------------
\newcommand{\errL}{\varepsilon_{2}}
\newcommand{\errI}{\varepsilon_{\infty}}
\newcommand{\errC}{\varepsilon_{2}^{\mathrm{c}}}
\newcommand{\errH}{\varepsilon_{H^{1}}}
\newcommand{\errT}{\varepsilon_{2}^{\mathrm{slice}}}

% ---- numbers ---------------------------------------------------------------------
\newcommand{\sci}[2]{#1\times 10^{#2}}

% ---- names -----------------------------------------------------------------------
\newcommand{\PH}{\textsf{Heat}}
\newcommand{\PWone}{\textsf{Wave1}}
\newcommand{\PWtwo}{\textsf{Wave2}}
\newcommand{\PLtwo}{\textsf{Lap2}}
\newcommand{\PLthree}{\textsf{Lap3}}
\newcommand{\PLthreeS}{\textsf{Lap3s}}
\newcommand{\Pone}{\textsf{LapD}}
\newcommand{\Ptwo}{\textsf{PoiD}}
\newcommand{\Pridge}{\textsf{RidgeD}}
\newcommand{\Pcospair}{\textsf{CosPairD}}
\newcommand{\Paltridge}{\textsf{AltRidgeD}}
\newcommand{\Psinlin}{\textsf{SinLinD}}
\newcommand{\Pexpcos}{\textsf{ExpCosD}}
\newcommand{\strat}[1]{\textsf{#1}}
\newcommand{\Adam}{Adam}
\newcommand{\LBFGS}{L-BFGS}
\newcommand{\AdamLBFGS}{Adam$\,\to\,$L-BFGS}
\newcommand{\code}[1]{\texttt{\detokenize{#1}}}

% ---- theorem environments (amsthm) -----------------------------------------------
\theoremstyle{plain}
\newtheorem{theorem}{Theorem}
\newtheorem{proposition}[theorem]{Proposition}
\newtheorem{lemma}[theorem]{Lemma}
\newtheorem{corollary}[theorem]{Corollary}
\theoremstyle{definition}
\newtheorem{definition}[theorem]{Definition}
\newtheorem{observation}[theorem]{Numerical observation}
\newtheorem{conjecture}[theorem]{Conjecture}
\theoremstyle{remark}
\newtheorem{remark}[theorem]{Remark}
